\begin{abstract}
    Il presente lavoro di tesi descrive la progettazione e lo sviluppo di un software full-stack
    a supporto del reparto commerciale di Cepi S.p.A.\cite{Cepi}, 
    contesto aziendale in cui l'autore opera come dipendente.
    L'azienda metalmeccanica forlivese è specializzata nella produzione di impianti
    di stoccaggio con logica \ac{ETO}.
    Il sistema realizzato si qualifica come soluzione \ac{CPQ},
    consentendo la configurazione del prodotto, la generazione di offerte
    commerciali e la conseguente stipula dei contratti.
    L'elaborato copre l'intero ciclo di vita del progetto: dall'analisi
    dell'as-is e la raccolta dei requisiti, fino alla modellazione del database relazionale e
    all'implementazione del software.
    La gestione dei dati è affidata a Microsoft SQL Server, mentre l'applicativo è stato sviluppato
    in linguaggio \csharp{} su framework \dotnet{}. La scelta di questo ecosistema risponde all'esigenza di
    realizzare un'applicazione cross-platform installabile in locale, requisito essenziale per garantire la piena
    operatività offline anche in contesti privi di connettività di rete.
    L'interfaccia utente è stata infine implementata sfruttando il framework \maui{}.
\end{abstract}